%% Updated main text section 2.3 — rank-based candidate gene scoring
\subsection{2.3 Candidate-gene prioritization}
To prioritize genes for downstream interpretation, CATFISH integrates three evidence
layers: locus-level GWAS support, MAGMA gene-level association, and CATFISH
pathway-level enrichment. For each gene $g$, we defined a composite score based on
rank-normalized evidence across all three layers,
\[
\mathrm{Score}(g)
= S_{\mathrm{GWAS}}(g)
+ S_{\mathrm{MAGMA}}(g)
+ 0.5\,S_{\mathrm{PATH}}(g)
+ 0.3\,B_{\mathrm{hub}}(g),
\]
where each rank score $S(\cdot) = -\log_{10}(\mathrm{rank}/N)$ places all evidence
sources on a common scale, and $B_{\mathrm{hub}}(g) = \log_{10}(n_{\mathrm{pathways}}(g)+1)$
rewards genes with membership in multiple pathways. Genes with no GWAS locus overlap
receive $S_{\mathrm{GWAS}}(g)=0$; genes not annotated to any pathway receive
$S_{\mathrm{PATH}}(g)=0$ and $B_{\mathrm{hub}}(g)=0$. Pathway evidence is downweighted
to half the strength of GWAS or MAGMA (weight $0.5$) because pathway $p$-values are
set-level statistics shared across all genes in a pathway rather than gene-specific
measurements. For interpretability, we also computed a discrete layer support count
$L_g = \mathbf{1}\{\mathrm{GWAS}\} + \mathbf{1}\{\mathrm{MAGMA}\} + \mathbf{1}\{\mathrm{PATH}\}$
and prioritized genes with $L_g \geq 2$ as multi-layer candidates for biological
follow-up. Full details of the scoring procedure are given in Appendix~A5.


%% =============================================================================
%% Updated appendix section A5 — rank-based candidate gene scoring
%% =============================================================================

\section{A5 Candidate-gene prioritization score}\label{app:candidate_score}

The candidate-gene prioritization score was designed as a heuristic ranking device to
identify genes supported across multiple analytical layers rather than as a formal
inferential model.

\paragraph{Rank normalization.}
For each gene $g$, we defined rank scores from three evidence layers. Let
$r_{\mathrm{GWAS}}(g)$ denote the rank of gene $g$ by its minimum within-locus GWAS
$p$-value among all genes with at least one overlapping SNP, $r_{\mathrm{MAGMA}}(g)$
the rank by MAGMA gene-level $p$-value among all genes, and $r_{\mathrm{PATH}}(g)$ the
rank of the gene's best-ranked pathway among all $N_{\mathrm{pathways}}$ tested pathways.
The corresponding rank scores are
\[
S_{\mathrm{GWAS}}(g)
  = -\log_{10}\!\left(\frac{r_{\mathrm{GWAS}}(g)}{N_{\mathrm{genes}}}\right), \qquad
S_{\mathrm{MAGMA}}(g)
  = -\log_{10}\!\left(\frac{r_{\mathrm{MAGMA}}(g)}{N_{\mathrm{genes}}}\right), \qquad
S_{\mathrm{PATH}}(g)
  = -\log_{10}\!\left(\frac{r_{\mathrm{PATH}}(g)}{N_{\mathrm{pathways}}}\right),
\]
where $N_{\mathrm{genes}}$ is the total number of genes with MAGMA results. Genes not
overlapping any GWAS locus receive $S_{\mathrm{GWAS}}(g)=0$; genes not annotated to
any pathway receive $S_{\mathrm{PATH}}(g)=0$.

\paragraph{Multi-pathway hub bonus.}
Genes annotated to multiple pathways receive an additional term
\[
B_{\mathrm{hub}}(g) = \log_{10}\!\left(n_{\mathrm{pathways}}(g)+1\right),
\]
where $n_{\mathrm{pathways}}(g)$ is the number of distinct pathways containing gene
$g$. This rewards biological connectivity without allowing multi-pathway membership to
dominate the score.

\paragraph{Composite score.}
The final prioritization score is
\[
\mathrm{Score}(g)
= S_{\mathrm{GWAS}}(g)
+ S_{\mathrm{MAGMA}}(g)
+ 0.5\,S_{\mathrm{PATH}}(g)
+ 0.3\,B_{\mathrm{hub}}(g).
\]

\paragraph{Rationale for weights.}
Because all rank scores share the same $-\log_{10}(\mathrm{rank}/N)$ scale, the weights
are directly interpretable as relative importance multipliers. Pathway evidence receives
weight $0.5$ (half that of GWAS or MAGMA) because a pathway $p$-value is a set-level
statistic distributed across all genes in that pathway rather than a gene-specific
measurement; awarding full weight would overstate gene-level resolution. GWAS and MAGMA
receive equal weight because they are regarded as complementary but equivalent sources
of direct genetic evidence at this stage of integration; any advantage of MAGMA in
aggregating multi-SNP signal is already captured implicitly through its typically more
significant gene-level $p$-values, which translate to higher rank scores naturally.
The hub bonus coefficient of $0.3$ is appreciable but subordinate: a gene in five
pathways gains $0.3\times\log_{10}(6)\approx 0.23$ additional score units, sufficient
to improve ranking within a tier but insufficient to promote a gene lacking primary
GWAS or MAGMA support.

\paragraph{Discrete layer support annotation.}
Separately from the composite score, we annotated each gene with binary support
indicators
\[
\mathbf{1}\{\mathrm{GWAS}\}, \quad
\mathbf{1}\{\mathrm{MAGMA}\}, \quad
\mathbf{1}\{\mathrm{PATH}\},
\]
where a gene was counted as a GWAS hit if its minimum locus $p$-value fell within the
top 1\% of all genes, a MAGMA hit if its MAGMA $p$-value fell within the top 5\%, and
a pathway hit if it belonged to at least one annotated pathway. The sum
$L_g = \mathbf{1}\{\mathrm{GWAS}\}+\mathbf{1}\{\mathrm{MAGMA}\}+\mathbf{1}\{\mathrm{PATH}\}$
was reported alongside the composite score; genes with $L_g \geq 2$ were designated
multi-layer candidates and prioritized for biological follow-up. The binary indicators
and composite score are consistent: a gene with higher $L_g$ will almost always
outrank a gene with lower $L_g$ because the rank-score contributions from additional
layers typically exceed the continuous within-layer differences.
